%% ma: L10-B15-0111
%% lop: 10
%% chuong: 6
%% bai: 15
%% dang: 10.15.6
%% loai: TLN
%% muc: TH
%% dap_an: "1"
%% nguon: "Ảnh giáo viên gửi 25/11/2026 (lô 2): Bài tập trắc nghiệm – Toán 10 Bài 15 Hàm số (sự biến thiên, đồ thị), Phần III Câu 3"
%% kien_thuc: "Lớp 10 Bài 15 (hàm số trong thực tế); bất phương trình bậc nhất"
%% trang_thai: chinh-thuc
%% ghi_chu: "Cùng tình huống với L10-B15-0086 (câu đó hỏi a+b+c+d), câu này hỏi lựa chọn nhà hàng."
\begin{ex}
Anh T cần đặt một tiệc cưới ước tính khoảng $30$ đến $35$ bàn. Nhà hàng $1$ đề nghị anh T đóng tiền cố định $20$ triệu đồng, sau khi tiệc cưới diễn ra sẽ đóng khoản còn lại với số tiền $2$ triệu đồng/1 bàn. Nhà hàng $2$ đề nghị anh T đóng tiền cố định $10$ triệu đồng, sau khi tiệc cưới diễn ra sẽ đóng khoản còn lại với số tiền $2{,}5$ triệu/1 bàn. Để tiết kiệm được chi phí cho tiệc cưới, anh T nên lựa chọn nhà hàng $n$ (giả sử rằng chất lượng phục vụ hai nhà hàng trên là ngang nhau)? Giá trị $n$ bằng bao nhiêu?
\shortans{1}
\loigiai{Với $x$ bàn: $f(x)=20+2x$, $g(x)=10+2{,}5x$. $f(x)-g(x)=10-0{,}5x<0$ với mọi $x\in[30;35]$. Vậy chọn nhà hàng $1$, $n=1$.}
\end{ex}
